% Part: sets-functions-relations
% Chapter: arithmetization
% Section: cauchy

\documentclass[../../../include/open-logic-section]{subfiles}

\begin{document}

\olfileid[fr]{sfr}{arith}{cauchy}

\olsection{Annexe : les réels construits à partir des suites de Cauchy}

Dans la \olref[cuts]{sec}, nous avons construit les réels comme coupures
de Dedekind. Nous allons présenter ici une autre construction. Elle
s'appuie sur la définition, due à Cauchy, de ce que nous appelons
aujourd'hui une suite de Cauchy. Son emploi pour \emph{construire} les
réels est toutefois dû à d'autres auteurs du dix-neuvième siècle,
notamment Weierstrass, Heine, M\'{e}ray et Cantor. Pour un exposé
historique, voir \citeauthor{OConnorRobertson:RN}
\citeyear{OConnorRobertson:RN}.

Avant d'aborder le dix-neuvième siècle, arrêtons-nous sur Simon Stevin
(1548--1620). Pour résumer, Stevin a compris que l'on pouvait envisager
chaque réel à partir de son développement décimal. Même un nombre
irrationnel tel que $\sqrt{2}$ possède un développement décimal dont
le début est :
\[
1{,}41421356237\ldots
\]
Il est très facile de représenter ces développements en théorie des
ensembles : il suffit de les considérer comme des fonctions
$d \colon \Nat \to \Nat$, où $d(n)$ est le chiffre décimal de rang $n$
qui nous intéresse. Il faut ensuite adapter un peu cette représentation
pour tenir compte de ce qui précède la virgule, ici simplement $1$.
Une autre adaptation, au moyen d'une relation d'équivalence, garantit
par exemple que $0{,}999\ldots=1$. Il n'est cependant pas difficile de
donner, dans l'esprit de Stevin, une construction parfaitement rigoureuse
des réels en théorie des ensembles.
\begin{editorial}
Les valeurs de $d$ sont des chiffres : elles doivent donc appartenir à
$\{0,\ldots,9\}$, restriction laissée implicite dans la source. Pour
représenter tous les réels, il faut également coder le signe et la partie
entière, puis identifier les deux écritures éventuelles d'un même nombre.
\end{editorial}

Stevin ne sera pas notre sujet principal ; pour en savoir davantage,
voir \citealt{KatzKatz2012}. Considérons plutôt une idée très proche.
Au lieu de traiter directement le développement décimal de $\sqrt{2}$,
on peut considérer une \emph{suite} d'approximations rationnelles de
plus en plus précises, obtenues en conservant davantage de chiffres :
\[
1;\ 1{,}4;\ 1{,}414;\ 1{,}4142;\ 1{,}41421;\ldots
\]
L'idée de représenter les réels par des « approximations de plus en
plus précises » conduit à une nouvelle série d'observations, analogues
à celles qui ont guidé la construction de $\Rat$ à partir de $\Int$,
ou de $\Int$ à partir de $\Nat$ :
\begin{enumerate}
\item Tout réel admet un développement décimal, éventuellement infini.
\item L'information contenue dans un tel développement peut tout aussi
bien être représentée par une suite de nombres rationnels.
\item Une suite de rationnels peut être considérée comme une fonction
de $\Nat$ dans $\Rat$ : il suffit de prendre pour $f(n)$ le rationnel
de rang $n$ de la suite.
\end{enumerate}
Bien entendu, une fonction \emph{quelconque} de $\Nat$ dans $\Rat$
ne nous donne pas nécessairement un réel. Considérons par exemple :
\[
f(n)=\begin{cases}
1 &\text{si }n\text{ est impair}\\
0 &\text{si }n\text{ est pair}
\end{cases}
\]
La difficulté est que la suite $0,1,0,1,0,1,0,\ldots$ ne semble
« se rapprocher » d'aucun réel. Pour retenir des suites qui se rapprochent
effectivement d'un réel, il faut donc une condition correspondant à
l'existence d'une \emph{limite}.

Nous avons déjà rencontré la notion de limite dans la
\olref[his][set][limits]{sec}. Nous ne pouvons pourtant pas reprendre
exactement la même définition. L'expression « $(\forall\epsilon>0)$ »
y quantifiait implicitement sur les réels, et nous considérions des
limites de fonctions définies sur les réels. Invoquer cette définition
supposerait donc déjà les réels, alors que nous cherchons précisément
à les \emph{construire}. Heureusement, une notion étroitement liée à
celle de limite permet de procéder autrement.
\begin{defn}\ollabel{def:CauchySequence}
Une fonction $f:\Nat\to\Rat$ est une \emph{suite de Cauchy} si et
seulement si, pour tout $\epsilon\in\Rat$ strictement positif,
$(\exists\ell\in\Nat)(\forall m,n>\ell)|f(m)-f(n)|<\epsilon$,
où $m$ et $n$ parcourent les naturels.
\end{defn}

L'idée reste celle d'une précision prescrite : pour un degré de
précision mesuré par $\epsilon$, on trouve un rang $\ell$ au-delà
duquel les termes sont suffisamment proches les uns des autres.
Notre suite $1$, $1{,}4$, $1{,}414$, $1{,}4142$, $1{,}41421$\ldots
illustre cette idée : pour approcher $\sqrt{2}$ avec une erreur
inférieure à $\nicefrac{1}{10^n}$, il suffit de considérer les termes
situés après celui de rang $n$.
\begin{editorial}
Le langage de la limite réelle reste ici une motivation intuitive,
ou renvoie au réel déjà construit par coupures. La définition de
Cauchy n'affirme pas qu'une limite existe dans $\Rat$ : elle compare
deux termes de la suite. Pour obtenir la condition de Cauchy à partir
d'approximations d'erreur donnée, on prend chaque erreur inférieure à
$\epsilon/2$ et on applique l'inégalité triangulaire.
\end{editorial}

On pourrait alors être tenté de dire qu'un réel \emph{est} simplement
une suite de Cauchy. Ce serait toutefois trop naïf, comme dans les
constructions de $\Int$ et de $\Rat$ : plusieurs suites de Cauchy
distinctes peuvent représenter le même réel. Pour le voir, prenons une
suite de Cauchy $f$ et définissons $g$ exactement comme $f$, sauf que
$g(0)\neq f(0)$. Les deux suites coïncident après le premier terme.
Le critère de la \olref{def:CauchySequence} est donc satisfait par l'une
comme par l'autre, et nous voudrons, au terme de la construction, leur
attribuer la même limite. Elles doivent ainsi « définir » le même réel
et être identifiées en tant que représentations de ce réel.

Il nous faut donc, une nouvelle fois, définir une relation d'équivalence
sur les suites de Cauchy et identifier les réels à des classes
d'équivalence. Définissons d'abord ce que signifie tendre vers $0$.
Pour une fonction $h:\Nat\to\Rat$, on dit que \emph{$h$ tend vers $0$}
si et seulement si, pour tout $\epsilon\in\Rat$ strictement positif,
$(\exists\ell\in\Nat)(\forall n>\ell)|h(n)|<\epsilon$, où $n\in\Nat$.
\footnote{Comparer avec la définition de
$\lim_{x\mathord{\rightarrow}\infty}f(x)=0$ dans la
\olref[his][set][limits]{sec}.}
Pour deux fonctions $f,g:\Nat\to\Rat$, posons en outre
$(f-g)(n)=f(n)-g(n)$. Définissons alors :
\[
f\Realequiv g\quad\text{si et seulement si }(f-g)\text{ tend vers }0.
\]
Il faut vérifier que $\Realequiv$ est une relation d'équivalence ;
c'est bien le cas. On peut alors définir les réels comme les classes,
pour $\Realequiv$, de toutes les suites de Cauchy de $\Nat$ dans $\Rat$.
\begin{editorial}
La source écrit d'abord « relations d'équivalence » là où il faut
« classes d'équivalence ». De même, ce sont les classes de suites qui
formeront un corps ordonné dans les énoncés qui suivent. Une suite
représentante et sa classe sont désormais distinguées partout.
\end{editorial}

\begin{prob}
Posons $f(n)=0$ pour tout $n$ et $g(n)=\frac{1}{(n+1)^2}$.
Montrer que ce sont deux suites de Cauchy et, plus précisément, que
toutes deux tendent vers $0$, de sorte que $f\sim_\Real g$.
\end{prob}

Comme d'habitude, nous noterons $\equivrep{f}{\Realequiv}$ la classe
d'équivalence ayant $f$ pour !!a{element}. Pour alléger la lecture,
nous omettrons désormais l'indice et écrirons simplement
$\equivrep{f}{}$. Pour tout $q\in\Rat$, nous posons aussi
$q_\Real=\equivrep{c_q}{}$, où $c_q$ est la fonction constante définie
par $c_q(n)=q$ pour tout $n\in\Nat$. Nous définissons alors les
relations et les opérations de base sur les réels, par exemple :
\begin{align*}
\equivrep{f}{}+\equivrep{g}{}&=\equivrep{(f+g)}{}\\
\equivrep{f}{}\times\equivrep{g}{}&=\equivrep{(f\times g)}{}
\end{align*}
où $(f+g)(n)=f(n)+g(n)$ et $(f\times g)(n)=f(n)\times g(n)$.
Il faut bien sûr vérifier que $(f+g)$, $(f-g)$ et $(f\times g)$ sont
des suites de Cauchy lorsque $f$ et $g$ en sont ; c'est effectivement
le cas, et nous vous laissons cette vérification.
\begin{editorial}
Il faut également vérifier que les opérations sur les classes ne
dépendent pas des représentants choisis. La stabilité des suites de
Cauchy sous les opérations ne suffit pas, à elle seule, à l'établir.
\end{editorial}

Définissons enfin l'ordre. La classe $\equivrep{f}{}$ est
\emph{strictement positive} si et seulement si
$\equivrep{f}{}\neq0_\Real$ et
$(\exists\ell\in\Nat)(\forall n>\ell)0<f(n)$.
On pose $\equivrep{f}{}<\equivrep{g}{}$ si et seulement si
$\equivrep{(g-f)}{}$ est strictement positive. Il faut vérifier que
cette définition ne dépend pas du choix de la fonction représentante
dans chaque classe d'équivalence.
% Variante inactive conservée de la source. Le cas d'égalité des classes
% est nécessaire : l'inégalité large terme à terme, seule, ne descend
% pas au quotient.
%\begin{align*}
%\equivrep{f}{}\leq\equivrep{g}{}\text{ ssi }&\equivrep{f}{}=\equivrep{g}{}\text{ ou }\\
%&( \exists \ell\in\Nat)(\forall n\in\Nat)(\ell<n\lif f(n)\leq g(n))
%\end{align*}
Une fois cette vérification faite, les propriétés algébriques attendues
s'établissent assez facilement. Autrement dit :
\begin{editorial}
Nous avons corrigé $0_\Rat$ en $0_\Real$ dans la définition de
positivité. La condition que la classe soit non nulle est essentielle :
une suite peut être strictement positive à chaque rang et tendre vers
zéro. Pour justifier l'ordre et les inverses, on utilise le fait
suivant : si une suite de Cauchy $u$ ne tend pas vers zéro, il existe
un rationnel $\delta>0$ tel qu'à partir d'un certain rang,
$|u(n)|>\delta$ et tous les $u(n)$ ont le même signe. En effet, il
existe $\epsilon_0>0$ rationnel tel que des termes arbitrairement
tardifs aient une valeur absolue au moins égale à $\epsilon_0$.
Au-delà d'un rang où les écarts entre termes sont inférieurs à
$\epsilon_0/2$, tous les termes sont donc du même signe et de valeur
absolue supérieure à $\epsilon_0/2$. Ce fait montre aussi que le signe
de la classe ne change pas lorsqu'on remplace $u$ par une suite
équivalente. Pour une classe non nulle, on peut prendre comme
représentante de l'inverse la suite égale à $1/u(n)$ à partir de ce
rang, et à $1$ avant celui-ci. La vérification qu'elle est de Cauchy
utilise
$|1/u(m)-1/u(n)|\leq |u(m)-u(n)|/\delta^2$.
\end{editorial}
\begin{thm}\ollabel{thm:cauchyorderedfield}
Les classes d'équivalence de suites de Cauchy forment un corps ordonné.
\end{thm}

\begin{proof}
En exercice.
\end{proof}

\begin{prob}
Démontrer que les classes d'équivalence de suites de Cauchy forment
un corps ordonné.
\end{prob}

La propriété de la borne supérieure est plus difficile à établir pour
les réels ainsi construits ; nous allons donc en donner la démonstration.

\begin{thm}
Tout ensemble non vide et majoré de classes d'équivalence de suites
de Cauchy admet une borne supérieure.
\end{thm}

\begin{proof}[Esquisse de démonstration]
Soit $\mathcal A$ un ensemble non vide et majoré de classes de suites
de Cauchy. Pour conserver les notations de la construction ci-dessous,
notons $S$ l'ensemble de toutes les suites $h$ telles que
$\equivrep{h}{}\in\mathcal A$. Les éléments de $S$ sont donc des suites,
et ceux de $\mathcal A$ des classes. Il existe $p\in\Rat$ tel que
$p_\Real$ majore $\mathcal A$. Choisissons $r\in S$ ; il existe
$q\in\Rat$ tel que $q_\Real<\equivrep{r}{}$. Ces bornes rationnelles
existent parce que toute suite de Cauchy rationnelle est bornée :
on applique cette observation à une représentante d'un majorant
de $\mathcal A$, puis à $r$. Si la borne supérieure de $\mathcal A$
existe, elle se trouve donc entre $q_\Real$ et $p_\Real$, bornes comprises.

Nous allons nous rapprocher de cette borne supérieure simultanément
par valeurs inférieures et supérieures. Définissons deux fonctions
$f,g:\Nat\to\Rat$, destinées à l'approcher respectivement par valeurs
supérieures et inférieures. Posons d'abord :
\begin{align*}
f(0)&=p\\
g(0)&=q
\end{align*}
Puis, en posant $a_n=\frac{f(n)+g(n)}2$, définissons :
\footnote{Il s'agit d'une définition par récurrence. Nous n'avons
cependant pas encore justifié l'emploi de telles définitions.}
\begin{align*}
f(n+1)&=\begin{cases}
a_n &\text{si }(\forall h\in S)\equivrep{h}{}\leq(a_n)_\Real\\
f(n)&\text{sinon}
\end{cases}\\
g(n+1)&=\begin{cases}
a_n &\text{si }(\exists h\in S)\equivrep{h}{}\geq(a_n)_\Real\\
g(n)&\text{sinon}
\end{cases}
\end{align*}
Les suites $f$ et $g$ sont de Cauchy. On peut le vérifier assez
facilement ; nous laissons le détail de cette vérification en exercice.
Voici les observations utiles : $f$ est décroissante, $g$ est
croissante, les intervalles $[g(n),f(n)]$ sont emboîtés et
\[
0\leq f(n+1)-g(n+1)\leq\frac{f(n)-g(n)}2.
\]
Leur longueur est donc au plus $(p-q)/2^n$. Deux termes suffisamment
tardifs de l'une ou l'autre suite appartiennent au même intervalle,
aussi petit que voulu. En particulier, $(f-g)$ tend vers $0$,
d'où $\equivrep{f}{}=\equivrep{g}{}$.
\emph{Précision éditoriale :} la source affirme que l'écart est divisé
exactement par deux à chaque étape. Avec ses deux conditions larges,
les deux premières branches peuvent s'appliquer simultanément, lorsque
le milieu est un maximum de $\mathcal A$ ; l'écart devient alors nul.
L'inégalité ci-dessus couvre aussi ce cas. Par construction, pour
chaque $n$, $(f(n))_\Real$ majore $\mathcal A$, et il existe
$h\in S$ tel que $(g(n))_\Real\leq\equivrep{h}{}$.

Montrons d'abord que $\equivrep{f}{}$ majore $\mathcal A$,
c'est-à-dire que
$(\forall h\in S)\equivrep{h}{}\leq\equivrep{f}{}$.
Nous utiliserons le \olref{thm:cauchyorderedfield} au cours du raisonnement.
Soit $h\in S$ et supposons, pour obtenir une contradiction, que
$\equivrep{f}{}<\equivrep{h}{}$, donc que
$0_\Real<\equivrep{(h-f)}{}$. La suite $f$ étant de Cauchy et
décroissante, il existe $n\in\Nat$ tel que
$\equivrep{(c_{f(n)}-f)}{}<\equivrep{(h-f)}{}$.
En effet, la positivité de la classe de $h-f$ donne un écart rationnel
$\delta>0$ tel que $h(m)-f(m)>\delta$ à partir d'un certain rang ;
on choisit $n$ assez grand pour que, pour tout $m$ assez grand,
$|f(n)-f(m)|<\delta/2$. La suite
$(h-f)-(c_{f(n)}-f)$ est alors, à partir d'un certain rang,
strictement supérieure à $\delta/2$.
Ainsi :
\[
(f(n))_\Real=\equivrep{c_{f(n)}}{}
<\equivrep{f}{}+\equivrep{(h-f)}{}=\equivrep{h}{},
\]
ce qui contredit $\equivrep{h}{}\leq(f(n))_\Real$, vraie par construction.

Montrons enfin que $\equivrep{f}{}=\equivrep{g}{}$ est le
\emph{plus petit} majorant de $\mathcal A$. Soit $j$ une suite de
Cauchy telle que $\equivrep{j}{}<\equivrep{g}{}$.
Comme ci-dessus, en utilisant cette fois la croissance de $g$,
il existe $n\in\Nat$ tel que $\equivrep{j}{}<(g(n))_\Real$ :
on prend un écart strictement positif pour $g-j$, puis un rang
où les écarts entre termes de $g$ lui sont inférieurs de moitié.
Or, par construction, il existe $h\in S$ tel que
$(g(n))_\Real\leq\equivrep{h}{}$. On a donc
$\equivrep{j}{}<\equivrep{h}{}$, et $\equivrep{j}{}$ ne majore pas
$\mathcal A$.
\end{proof}

\end{document}
